\documentclass[11pt]{article}

\usepackage[margin=2.5cm]{geometry}
\usepackage{amsmath,amssymb}
\usepackage{booktabs}
\usepackage{longtable}
\usepackage{siunitx}
\usepackage[hidelinks]{hyperref}

\sisetup{per-mode=symbol}

\newcommand{\dd}{\,\mathrm{d}}
\newcommand{\mac}[1]{\left\langle #1 \right\rangle}
\newcommand{\inv}[1]{\left| #1 \right|}
\newcommand{\Tk}{\tilde{T}}
\newcommand{\NN}{\mathcal{N}}

\title{A multi-plate, depth-integrated elasto-viscoplastic model\\
of the lithosphere with emergent plate boundaries}
\author{}
\date{}

\begin{document}
\maketitle

\begin{abstract}
We describe a reduced model of plate tectonics on a sphere. Each tectonic plate
is a separate triangulated surface mesh whose vertices carry vertical columns of
lithosphere. A column is described by a small state: crustal volume, reference
densities of crust and lithospheric mantle (used as a proxy for rock type), the
thermal thickness of the plate, a hydration index, a damage index, a
depth-integrated stress resultant, and a flag recording whether the column has
been subducted. From this state, depth profiles of temperature, density, and
strength are reconstructed and integrated to give a constitutive law for each
column: a Maxwell elastic element in series with a Herschel--Bulkley
viscoplastic element, weakened by a thixotropic damage variable. Plates interact
only through their boundaries, which are not conforming: plates may separate,
overlap, or slide past one another. Spreading, subduction, transform faulting,
collision, and rifting are not prescribed as boundary labels but follow from
the local geometry of gaps, overlaps, and sliding, together with the state of
the columns on either side. Subducted columns are retained beneath the
overriding plate as a slab, from which slab pull, interface friction, and arc
processes are computed. Convection in the underlying mantle is supplied by a
separate model as a velocity field and an upwelling rate at the base of the
plates.
\end{abstract}

\tableofcontents

%======================================================================
\section{Overview}
%======================================================================

The lithosphere is treated as a collection of thin sheets in the sense of
England and McKenzie~\cite{EnglandMcKenzie1982}: horizontal length scales are
large compared with plate thickness, horizontal velocity is uniform with depth
within a column, vertical normal stress is lithostatic, and mechanical response
is described by depth-integrated stress resultants (force per unit length). A
plate is defined kinematically as the material that moves coherently with the
surface; its base is identified with a characteristic isotherm $T_L$.

The model has five layers of description:
\begin{enumerate}
  \item \emph{State}, stored per column (Section~\ref{sec:state}).
  \item \emph{Reconstructed depth profiles} of temperature, density, and
        strength, computed from the state through closure relations
        (Sections~\ref{sec:thermal}--\ref{sec:strength}).
  \item \emph{Depth-integrated column properties}
        (Section~\ref{sec:integrated}).
  \item \emph{Plate-interior field equations}: constitutive law, damage
        evolution, and force balance
        (Sections~\ref{sec:constitutive}, \ref{sec:damage},
        \ref{sec:forcebalance}).
  \item \emph{Boundary interactions} between plates, including the creation and
        destruction of plate area and changes in plate topology
        (Section~\ref{sec:boundaries}).
\end{enumerate}
Quantities in layers 2 and 3, and all properties of plate boundaries, are
recomputed from the state when needed rather than stored.

%======================================================================
\section{Geometry, coordinates, and notation}
%======================================================================

\paragraph{Plates and columns.}
The surface of a sphere of radius $R_\oplus$ is covered by a set of plates
$\{\Omega_p\}$. Each plate is a triangulated mesh. Each vertex $v$ of a plate
mesh represents a vertical column of lithosphere occupying the vertex's dual
cell, of area $A_v$ (one third of the area of each incident triangle).
Triangles are used to evaluate horizontal gradients; columns carry the state.
Each plate's mesh has one or more boundary loops, which are ordered polylines of
boundary vertices.

\paragraph{Coordinates.}
At each point, $x_i$ ($i\in\{1,2\}$) are Cartesian coordinates in the local
tangent plane, with summation over repeated horizontal indices. Two vertical
coordinates are used: $z$, depth below sea level, and $\zeta$, depth below the
local solid surface, both positive downward. If $e$ is the elevation of the
solid surface above sea level, then $z = \zeta - e$. On a plate boundary, $n_i$
is the outward unit normal of the plate in the tangent plane and $s_i$ the unit
tangent.

\paragraph{Fields and operators.}
Temperatures written $T$ are in degrees Celsius; the same temperature in kelvin
is $\Tk = T + 273.15$. Horizontal velocity is $u_i$, and
$\mathrm{D}/\mathrm{D}t = \partial_t + u_j\partial_j$ is the material
derivative. The strain-rate and spin tensors are
\begin{equation}
  \dot\varepsilon_{ij} = \tfrac12\left(\partial_i u_j + \partial_j u_i\right),
  \qquad
  W_{ij} = \tfrac12\left(\partial_i u_j - \partial_j u_i\right).
\end{equation}
For a symmetric horizontal tensor $A_{ij}$,
\begin{equation}
  \inv{A} = \left(A_{11}^2 + A_{22}^2 + A_{11}A_{22} + A_{12}^2\right)^{1/2}
  \label{eq:invariant}
\end{equation}
is the second invariant $\left(\tfrac12 a_{pq}a_{pq}\right)^{1/2}$ of the
traceless three-dimensional tensor whose horizontal block is $A$ and whose
vertical component is $-A_{kk}$; it is the natural invariant for an
incompressible thin sheet. We write $\dot\varepsilon_{II} =
\inv{\dot\varepsilon}$. The Macaulay bracket is $\mac{x} = \max(x,0)$, and
$\operatorname{clamp}(x,a,b) = \min(\max(x,a),b)$.

\paragraph{Regularized direction.}
For a velocity vector $a_i$ we write
$\hat a_i = a_i / \sqrt{a_k a_k + v_{\mathrm{reg}}^2}$, a unit vector
regularized at small speeds by $v_{\mathrm{reg}}$, used wherever frictional
forces oppose relative motion.

%======================================================================
\section{State}
\label{sec:state}
%======================================================================

Each column carries the state in Table~\ref{tab:state}. Membership of a column
in a plate is given by the mesh it belongs to.

\begin{table}[h]
\centering
\caption{Per-column state.}
\label{tab:state}
\begin{tabular}{@{}lp{7.2cm}lp{3.0cm}@{}}
\toprule
Symbol & Meaning & Units & Character \\
\midrule
$V_c$ & crustal volume of the column & \si{\cubic\metre} & conserved, transferable \\
$\rho_{0,c}$ & reference density of crust at $T_0$, $P=0$ (rock-type proxy) & \si{\kilo\gram\per\cubic\metre} & material attribute \\
$\rho_{0,m}$ & reference density of lithospheric mantle at $T_0$, $P=0$ & \si{\kilo\gram\per\cubic\metre} & material attribute \\
$L$ & thermal thickness of the plate (depth of the $T_L$ isotherm) & \si{\metre} & strained, thermally evolving \\
$w$ & hydration index, $0 \le w \le 1$ & -- & advected, with sources \\
$D$ & damage index, $0 \le D \le 1$ & -- & evolving \\
$T_{ij}$ & depth-integrated deviatoric stress resultant & \si{\newton\per\metre} & evolving \\
$\varsigma$ & subduction flag, $\varsigma\in\{0,1\}$ & -- & set once \\
\bottomrule
\end{tabular}
\end{table}

Crustal thickness is derived, $h_c = V_c / A_v$, so that horizontal shortening
of a column thickens its crust without a separate rule, and crustal mass is
conserved by construction. The plate is divided into two layers indexed by
$k\in\{c,m\}$: crust, $0 \le \zeta < h_c$, and lithospheric mantle,
$h_c \le \zeta < L$, with $h_c < L$.

The subduction flag is the one piece of state that exists only because plates
overlap. A column with $\varsigma = 1$ has passed beneath another plate at a
trench; it no longer takes part in the surface force balance but is retained as
part of a slab (Section~\ref{sec:convergent}). Its other state variables are
frozen at their values on subduction. The flag records which of two overlapping
plates is underneath, which cannot be recovered from geometry alone.

\paragraph{Mixing rule.}
When crust is transferred between columns (by accretion at a trench, arc
magmatism, or conservative remapping after remeshing), volumes add and the
reference density of the receiving column becomes the volume-weighted mean of
the combined crust. Intensive mantle attributes are averaged weighted by column
area.

%======================================================================
\section{Thermal structure}
\label{sec:thermal}
%======================================================================

Within a column the temperature is linear in depth:
\begin{equation}
  T(\zeta) = T_s + G\,\zeta, \qquad
  G = \frac{T_L - T_s}{L}, \qquad 0 \le \zeta \le L,
  \label{eq:geotherm}
\end{equation}
where $T_s$ is the surface temperature and $G$ the thermal gradient. Because
$T_s$ and $T_L$ are global constants, the thermal state of a column is fully
described by $L$. Below the plate, the convecting mantle follows an adiabat
$T_a(z) = T_p + \gamma_{\mathrm{ad}}\,z$ with potential temperature $T_p$ and
adiabatic gradient $\gamma_{\mathrm{ad}}$. The small mismatch between $T_L$ and
$T_a$ at the plate base is neglected.

%======================================================================
\section{Rock-property closure}
\label{sec:closure}
%======================================================================

Each layer $k$ has a brittle--ductile transition temperature $T_{bd,k}$ and a
creep activation energy $Q_k$, obtained by piecewise-linear interpolation in
$\rho_{0,k}$ through anchor points $\{(\rho^{(a)}, T_{bd}^{(a)}, Q^{(a)})\}$
(Table~\ref{tab:anchors}), with constant extrapolation outside the anchored
range. With interpolants $\mathcal{T}(\rho_0)$ and $\mathcal{Q}(\rho_0)$,
\begin{equation}
  T_{bd,k} = \mathcal{T}(\rho_{0,k}) - \Delta T_w\, w,
  \qquad
  Q_k = \mathcal{Q}(\rho_{0,k}).
  \label{eq:Tbd}
\end{equation}
Hydration lowers the transition temperature by up to $\Delta T_w$, and raises the
pore-fluid pressure ratio (pore pressure divided by lithostatic pressure):
\begin{equation}
  \lambda(w) = \lambda_0 + (\lambda_1 - \lambda_0)\, w,
  \label{eq:lambda}
\end{equation}
where $\lambda_0$ corresponds to a hydrostatic fluid column and $\lambda_1$ to a
strongly overpressured one.

\begin{table}[h]
\centering
\caption{Anchor points for the rock-property closure. Transition temperatures
are for the reference strain rate $\dot\varepsilon_{\mathrm{ref}}$ under
compression. Activation energies follow laboratory flow laws for wet quartzite
\cite{GleasonTullis1995}, dry diabase \cite{Mackwell1998}, and dry olivine
\cite{HirthKohlstedt2003}.}
\label{tab:anchors}
\begin{tabular}{@{}lccc@{}}
\toprule
Representative rock & $\rho^{(a)}$ (\si{\kilo\gram\per\cubic\metre}) & $T_{bd}^{(a)}$ (\si{\celsius}) & $Q^{(a)}$ (\si{\kilo\joule\per\mole}) \\
\midrule
felsic (quartz-dominated) & 2700 & 300 & 223 \\
mafic (diabase, granulite) & 2950 & 600 & 485 \\
ultramafic (olivine) & 3250 & 750 & 530 \\
\bottomrule
\end{tabular}
\end{table}

Within the mantle, compositional depletion lowers $\rho_{0,m}$ without weakening
the rock; the strength of depleted mantle is represented through low hydration.
Densities between the mafic and ultramafic anchors have little physical meaning
and should not arise in practice.

%======================================================================
\section{Density, buoyancy, isostasy, and gravitational potential energy}
\label{sec:density}
%======================================================================

Density within layer $k$ of a column and within the asthenosphere are
\begin{equation}
  \rho_k(\zeta) = \rho_{0,k}\left[1 - \alpha\left(T(\zeta) - T_0\right)\right],
  \qquad
  \rho_a(z) = \rho_{a,0}\left[1 - \alpha\left(T_a(z) - T_0\right)\right],
\end{equation}
where $\alpha$ is the thermal expansion coefficient, $T_0$ the reference
temperature at which reference densities are specified, and $\rho_{a,0}$ the
reference density of undepleted asthenospheric mantle. Compressibility is
neglected, since it is nearly common to all columns.

\paragraph{Column buoyancy.}
The net buoyancy of a column per unit area, relative to asthenosphere at the
same depth below its surface, is
\begin{equation}
  B = \int_0^{L} \left[\rho_a(\zeta) - \rho(\zeta)\right] \dd\zeta,
  \label{eq:buoyancy}
\end{equation}
where $\rho(\zeta)$ is $\rho_c$ or $\rho_m$ according to layer. Positive $B$
indicates a column lighter than the mantle it displaces; continental columns
have $B > 0$ and old oceanic columns $B < 0$. Buoyancy decides which plate
subducts at a convergent boundary and when subduction must stop.

\paragraph{Isostatic elevation.}
For a compensation depth $z_c$ below the deepest plate base, the mass per unit
area above $z_c$ for a column with surface elevation $e$ is
\begin{equation}
  M(e) = \rho_w \mac{-e}
  + \int_{-e}^{-e+h_c} \rho_c \dd z
  + \int_{-e+h_c}^{-e+L} \rho_m \dd z
  + \int_{-e+L}^{z_c} \rho_a(z) \dd z,
\end{equation}
where $\rho_w$ is seawater density. Local isostasy requires
$M(e) = M_{\mathrm{ref}}$, with $M_{\mathrm{ref}}$ fixed by a reference column
of known elevation $e_{\mathrm{ref}}$. Since $\mathrm{d}M/\mathrm{d}e > 0$, the
root is unique.

\paragraph{Pressure and gravitational potential energy.}
With $z_{\mathrm{top}} = \min(0,-e)$,
\begin{equation}
  P(z) = g \int_{z_{\mathrm{top}}}^{z} \rho(z')\dd z',
  \qquad
  \Gamma = \int_{z_{\mathrm{top}}}^{z_c} P(z)\dd z - \Gamma_{\mathrm{ref}},
  \label{eq:GPE}
\end{equation}
where $g$ is gravitational acceleration and $\Gamma$ is the gravitational
potential energy per unit area, measured relative to that of the reference
column, $\Gamma_{\mathrm{ref}}$. Subtracting the constant reduces round-off,
since only differences of $\Gamma$ affect the dynamics. We write
$P_s(\zeta) = P(\zeta - e)$ for pressure as a function of depth below the solid
surface. $\Gamma_0$ denotes the value of $\Gamma$ for a newly formed column
(Section~\ref{sec:spreading}).

%======================================================================
\section{Strength profile}
\label{sec:strength}
%======================================================================

The strength of a column at depth $\zeta$ is the differential stress
$\sigma = \sigma_1 - \sigma_3$ it can sustain at the reference strain rate
$\dot\varepsilon_{\mathrm{ref}}$.

\paragraph{Frictional strength.}
Frictional sliding on optimally oriented faults with friction coefficient
$\mu_f$ \cite{Byerlee1978} limits the differential stress to
\begin{equation}
  \sigma_b(\zeta;\chi) = \beta(\chi)\,\big(1-\lambda(w)\big)\,P_s(\zeta).
  \label{eq:brittle}
\end{equation}
With $R = (\sqrt{1+\mu_f^2}+\mu_f)^2$, the coefficients for thrust, strike-slip,
and normal faulting are
\begin{equation}
  \beta_c = R - 1, \qquad
  \beta_s = \frac{2(R-1)}{R+1}, \qquad
  \beta_e = 1 - \frac{1}{R},
\end{equation}
approximately $2.12$, $1.03$, and $0.68$ for $\mu_f = 0.6$. The deformation mode
is read from the stress resultant,
\begin{equation}
  \chi = \frac{T_{kk}}{\sqrt3\,\inv{T}} \in [-1,1]
  \qquad (\chi = 0 \text{ if } \inv{T}=0),
  \label{eq:chi}
\end{equation}
negative in horizontal compression and positive in horizontal tension (tension
positive), and
\begin{equation}
  \beta(\chi) =
  \begin{cases}
    \beta_s + (\beta_s - \beta_c)\,\chi, & \chi < 0,\\
    \beta_s + (\beta_e - \beta_s)\,\chi, & \chi \ge 0.
  \end{cases}
\end{equation}

\paragraph{Creep strength.}
Power-law creep, $\dot\varepsilon = A\,\sigma^n\exp(-Q/R_g\Tk)$, is anchored at
the depth of the brittle--ductile transition of layer $k$,
$\zeta_{bd,k} = (T_{bd,k} - T_s)/G$, where it is required to equal the
compressional frictional strength under hydrostatic pore pressure:
\begin{equation}
  \sigma_{d,k}(\zeta) = \sigma^*_k
  \exp\!\left[\frac{Q_k}{n R_g}\left(\frac{1}{\Tk(\zeta)} - \frac{1}{\Tk_{bd,k}}\right)\right],
  \qquad
  \sigma^*_k = \beta_c(1-\lambda_0)\,P_s(\zeta_{bd,k}),
  \label{eq:ductile}
\end{equation}
with stress exponent $n$ and gas constant $R_g$. This replaces the
pre-exponential factor $A$ by the more intuitive $T_{bd,k}$. The anchor uses
$\lambda_0$ so that pore pressure acts only on frictional strength. If
$\zeta_{bd,k}$ lies outside the layer it still serves as the anchor. At other
strain rates, creep strength scales as
$(\dot\varepsilon/\dot\varepsilon_{\mathrm{ref}})^{1/n}$.

\paragraph{Low-temperature plasticity and damage.}
A cap $\sigma_{\max}$ represents low-temperature plasticity. Damage weakens the
frictional and plastic limbs by $(1-a_bD)$ and the creep limb by $(1-a_dD)$.

\paragraph{Composite profile.}
For $\zeta$ in layer $k$,
\begin{equation}
  \sigma(\zeta) = \min\Big\{
    (1-a_bD)\,\sigma_b(\zeta;\chi),\;
    (1-a_dD)\,\sigma_{d,k}(\zeta),\;
    (1-a_bD)\,\sigma_{\max}
  \Big\}.
  \label{eq:envelope}
\end{equation}
Let $\mathcal{B}$ be the set of depths in $[0,L)$ where the frictional or plastic
limb is the minimum, and $\mathcal{V}$ where the creep limb is.

%======================================================================
\section{Depth-integrated column properties}
\label{sec:integrated}
%======================================================================

\paragraph{Yield and creep resultants.}
\begin{equation}
  F_b(\chi) = \kappa_\sigma \int_{\mathcal{B}} \sigma\dd\zeta,
  \qquad
  F_d(\chi) = \kappa_\sigma \int_{\mathcal{V}} \sigma\dd\zeta,
  \label{eq:resultants}
\end{equation}
where $\kappa_\sigma \in [\tfrac12, 1/\sqrt3]$ converts differential stress to
the invariant of Eq.~\eqref{eq:invariant}. $F_b$ is the rate-independent yield
resultant and $F_d$ the creep resultant at $\dot\varepsilon_{\mathrm{ref}}$. The
integrals are evaluated by quadrature.

\paragraph{Elastic thickness and flexural rigidity.}
Let $\mathcal{S} = \{\zeta\in[0,L): \sigma(\zeta)\ge\sigma_{\mathrm{thr}}\}$, and
let $\ell_1,\dots,\ell_q$ be the lengths of its maximal connected intervals.
Following Burov and Diament~\cite{BurovDiament1995},
\begin{equation}
  H_e = \sum_p \ell_p,
  \qquad
  T_e = \Big(\sum_p \ell_p^3\Big)^{1/3},
  \qquad
  D_f = \frac{E\,T_e^3}{12(1-\nu^2)},
\end{equation}
the membrane elastic thickness, the effective elastic thickness for bending, and
the flexural rigidity, with Young's modulus $E$ and Poisson's ratio $\nu$.
Mechanically separated strong layers bend independently and combine as a sum of
cubes. $D_f$ is needed only if flexure is included; otherwise vertical position
follows local isostasy.

\paragraph{Closed-form approximation.}
For one layer of uniform density $\bar\rho$ with no damage, no cap, and thrust
faulting, let $k_B = \beta_c(1-\lambda)\bar\rho g$ and linearize the Arrhenius
term about $\Tk_{bd}$. Creep strength then decays below $\zeta_{bd}$ with
e-folding depth $\delta = nR_g\Tk_{bd}^{\,2}/(Q\,G)$, and
\begin{equation}
  F_b \approx \tfrac12\kappa_\sigma k_B\zeta_{bd}^2,
  \qquad
  F_d \approx \kappa_\sigma k_B\zeta_{bd}\,\delta
  \left(1 - e^{-(L-\zeta_{bd})/\delta}\right),
  \qquad
  H_e \approx \zeta_{bd} + \delta\ln\frac{k_B\zeta_{bd}}{\sigma_{\mathrm{thr}}}.
\end{equation}
For mantle with $G\approx\SI{13}{\kelvin\per\kilo\metre}$, $\delta\approx
\SI{4}{\kilo\metre}$; for quartz-rich crust with
$G\approx\SI{18}{\kelvin\per\kilo\metre}$, $\delta\approx\SI{3}{\kilo\metre}$.
These expressions serve as checks on the quadrature.

%======================================================================
\section{Constitutive law}
\label{sec:constitutive}
%======================================================================

The stress resultant is $T_{ij} = \int_0^L \tau_{ij}\dd\zeta$, where $\tau_{ij}$
is the horizontal block of the deviatoric stress. Each column behaves as an
elastic element in series with a viscoplastic element:
\begin{equation}
  \dot\varepsilon_{ij}
  = \frac{1}{2G_sH_e}\frac{\mathcal{D}T_{ij}}{\mathcal{D}t}
  + \Phi\!\left(\inv{T}\right)\frac{T_{ij}}{\inv{T}},
  \qquad
  \Phi\!\left(\inv{T}\right) = \dot\varepsilon_{\mathrm{ref}}
  \mac{\frac{\inv{T} - F_b(\chi)}{F_d(\chi)}}^{n},
  \label{eq:evp}
\end{equation}
where $G_s$ is the elastic shear modulus and
\begin{equation}
  \frac{\mathcal{D}T_{ij}}{\mathcal{D}t}
  = \frac{\mathrm{D}T_{ij}}{\mathrm{D}t} - W_{ik}T_{kj} + T_{ik}W_{kj}
\end{equation}
is the co-rotational (Jaumann) derivative; on the sphere, $T_{ij}$ is also
parallel-transported between tangent planes. The viscoplastic strain rate has
invariant $\Phi$. Below the yield resultant the column is elastic; above it, it
flows as a depth-integrated Herschel--Bulkley material. This is an
elasto-viscoplastic law of the type of Saramito~\cite{Saramito2009}. For
implementation the viscoplastic part may be written
$T_{ij} = 2\bar\eta\,\dot\varepsilon^{vp}_{ij}$ with
$\bar\eta = \inv{T}/(2\Phi)$, regularized near yield by a small background
viscosity.

The limit $F_b\to0$, $H_e\to\infty$ is the power-law thin viscous sheet
\cite{EnglandMcKenzie1982}; the limit $F_b, F_d, H_e \to \infty$ is a rigid
plate (Section~\ref{sec:solution}). The partition into $\mathcal{B}$ and
$\mathcal{V}$ is evaluated at $\dot\varepsilon_{\mathrm{ref}}$ and held fixed,
an approximation justified by the weak dependence of creep strength on strain
rate ($\propto\dot\varepsilon^{1/n}$).

%======================================================================
\section{Damage}
\label{sec:damage}
%======================================================================

Damage is a structural variable that is broken down by deformation and rebuilt
by healing, as in models of thixotropic fluids; geologically it represents fault
weakening and grain-size reduction \cite{BercoviciRicard2012}. For unsubducted
columns,
\begin{equation}
  \frac{\mathrm{D}D}{\mathrm{D}t}
  = k_D\,\Phi\,(1-D)
  - \frac{D}{\tau_h}
  + \partial_i\!\left(\ell_D^2\,k_D\,\Phi\,\partial_i D\right)
  + S_D^{\Upsilon},
  \label{eq:damage}
\end{equation}
with $D$ clamped to $[0,1]$. The terms are, in order: breakdown in proportion to
viscoplastic strain rate, with $k_D$ the damage per unit strain; healing with
time scale
\begin{equation}
  \tau_h = \tau_{h,0}\exp\!\left[\frac{Q_h}{R_g}\left(\frac{1}{\Tk_h}-\frac{1}{\Tk_{h,0}}\right)\right],
  \qquad T_h = T_s + G\,\zeta_h,
\end{equation}
evaluated at a representative depth $\zeta_h$ in the strong part of the plate;
diffusion active only where the column is deforming, which sets the width of
shear zones to the physical length $\ell_D$ rather than the mesh spacing and
prevents damage in inactive regions from spreading; and a source from mantle
upwelling, defined in Section~\ref{sec:rifting}. Damage diffusion acts within a
plate and has no flux across plate boundaries.

%======================================================================
\section{Plate boundaries}
\label{sec:boundaries}
%======================================================================

Plate boundaries are not conforming. Each plate's boundary vertices move with
that plate, so neighbouring plates may separate (leaving a gap), overlap, or
slide. The kind of boundary at any point is not a stored label; it follows from
the local geometry and from the state of the columns on either side, as follows.

%----------------------------------------------------------------------
\subsection{Contact geometry}
%----------------------------------------------------------------------

For each boundary vertex $b$ of plate $p$, the nearest boundary segment of any
other plate $q$ is found with a global spatial index. The signed gap $g_b$ is
the distance from $b$ to that segment, positive if $b$ lies outside $\Omega_q$
and negative if inside. The relative velocity is
$\Delta u_i = u_i^{(p)} - u_i^{(q)}$ evaluated at $b$ and at its projection onto
$q$, with normal and tangential components
\begin{equation}
  v_n = \Delta u_i n_i, \qquad v_t = \Delta u_i s_i,
\end{equation}
where $n_i$ is the outward normal of $p$. Thus $v_n > 0$ is convergence and
$v_n < 0$ divergence. Boundary vertices are said to be in contact when
$|g_b| \le g_{\mathrm{tol}}$.

%----------------------------------------------------------------------
\subsection{Classification}
\label{sec:classification}
%----------------------------------------------------------------------

At each boundary point exactly one of the following holds:
\begin{enumerate}
  \item \textbf{Gap} ($g_b > g_{\mathrm{tol}}$): the boundary is free, and new
        plate is created in the gap (Section~\ref{sec:spreading}). This is a
        spreading boundary.
  \item \textbf{Overlap with slab}: subducted columns of $q$ lie beneath $p$ at
        this point, or vice versa. This is a subduction boundary
        (Section~\ref{sec:convergent}).
  \item \textbf{Contact or overlap without slab}: the plates are in contact or
        have begun to overlap without either having subducted. The boundary
        transmits normal force and frictional shear
        (Section~\ref{sec:contact}), and may initiate subduction or weld the
        plates together.
\end{enumerate}
A contact boundary dominated by tangential motion is a transform fault; one
dominated by convergence is a nascent subduction zone or a collision. The memory
that would otherwise require stored boundary labels is carried by the geometry
(gaps filled by new plate) and by the retained slab (subducted columns persist
until they sink out).

%----------------------------------------------------------------------
\subsection{Spreading boundaries}
\label{sec:spreading}
%----------------------------------------------------------------------

Where the gap between plates $p$ and $q$ exceeds $\Delta x_{\mathrm{acc}}$, new
plate is created. Vertices are inserted along the midline of the gap,
duplicated so that each plate receives its own copy, and each plate's mesh is
extended by a strip of new triangles between its former boundary and the
midline. New columns are initialized as
\begin{equation}
  L = L_0,\quad
  V_c = h_c^{o}A_v,\quad
  \rho_{0,c} = \rho_{0,c}^{o},\quad
  \rho_{0,m} = \rho_{0,m}^{o},\quad
  w = w^{o},\quad
  D = 0,\quad
  T_{ij} = 0,\quad
  \varsigma = 0.
\end{equation}
Boundary edges longer than $\ell_{\max}$ are split. The location and rate of
spreading are not prescribed: plates separate wherever the force balance
(Section~\ref{sec:forcebalance}) drives them apart, and new plate fills the
space. The gravitational potential energy difference between young and old
lithosphere (``ridge push'') enters through $\Gamma$; no additional ridge force
is applied. Gaps arising anywhere, not only at established ridges, are filled
the same way.

%----------------------------------------------------------------------
\subsection{Contact boundaries}
\label{sec:contact}
%----------------------------------------------------------------------

\paragraph{Normal force.}
Interpenetration without subduction is resisted by a normal contact force per
unit length,
\begin{equation}
  N_c = k_c \mac{-g_b},
\end{equation}
with penalty stiffness $k_c$ (or, equivalently, a Lagrange multiplier enforcing
$g_b \ge 0$). It acts on $p$ as $-N_c n_i$ and equally and oppositely on $q$.

\paragraph{Frictional shear.}
Sliding is resisted by the strike-slip yield resultant of the weaker of the two
boundary columns, which already includes the effects of hydration and damage:
\begin{equation}
  F_T = \min\left(F_b^{(p)}(0),\,F_b^{(q)}(0)\right),
  \qquad
  f^{T}_i = -F_T\,\widehat{(v_t s)}_i \quad \text{on } p,
\end{equation}
and oppositely on $q$. No separate fault friction parameter is introduced.

\paragraph{Subduction initiation.}
If the normal force exceeds the compressional yield resultant of the boundary
column with the lower buoyancy $B$, and that buoyancy is below the threshold for
subduction,
\begin{equation}
  N_c > F_b^{(\mathrm{dense})}(-1)
  \quad\text{and}\quad
  B^{(\mathrm{dense})} < B_{\mathrm{crit}},
\end{equation}
the denser plate begins to subduct along the contiguous stretch of boundary
where these conditions hold: its columns are flagged $\varsigma=1$ as they pass
beneath the other plate (Section~\ref{sec:convergent}).

\paragraph{Welding.}
If both boundary columns have $B \ge B_{\mathrm{crit}}$ and the plates are
converging, neither can subduct. The two meshes are stitched along the contact
into a single plate, and the stitched columns are given $D = 1$. The resulting
suture is a weak zone that heals according to Eq.~\eqref{eq:damage} and may be
reactivated later.

%----------------------------------------------------------------------
\subsection{Subduction boundaries}
\label{sec:convergent}
%----------------------------------------------------------------------

Let plate $p$ be subducting beneath plate $q$. The trench is the part of
$q$'s boundary that overlies subducted columns of $p$.

\paragraph{Subduction of columns.}
A column of $p$ with $\varsigma = 0$ is flagged $\varsigma = 1$ when it passes
beneath $q$ by more than $g_{\mathrm{tol}}$. At that moment its crust is
partitioned according to its rock type. If $\rho_{0,c} < \rho_{\mathrm{scr}}$
(buoyant, felsic crust), the entire crustal volume is scraped off: a fraction
$f_{\mathrm{acc}}$ is transferred to the nearest trench vertex of $q$, and the
remainder to the nearest unsubducted boundary vertex of $p$, using the mixing
rule. Denser crust is retained in the slab.

\paragraph{Slab columns.}
Subducted columns keep their mesh connectivity and frozen state, take no part
in the force balance of Eq.~\eqref{eq:balance}, and move with the velocity of
the nearest unsubducted boundary vertex of $p$. The downdip length of slab
behind each subducted column is not stored but recomputed as its map distance
$d_v$ from the trench, which grows at the convergence rate regardless of trench
migration. A subducted column is deleted when
\begin{equation}
  d_v > d_{\max} = \frac{z_{\mathrm{slab}}}{\sin\theta_d},
\end{equation}
where $\theta_d$ is a representative slab dip and $z_{\mathrm{slab}}$ the depth
below which the slab no longer pulls on the plate (for example, where it is
supported at the base of the upper mantle).

\paragraph{Slab pull.}
Each subducted column has a negative buoyancy computed from its frozen state by
Eq.~\eqref{eq:buoyancy}, with its retained crust converted to eclogite (density
$\rho_{\mathrm{ecl}}$ in place of $\rho_{0,c}$). Its downdip weight is
$W_v = g\sin\theta_d\,\mac{-B_v^{\mathrm{slab}}}\,A_v$. Each subducted column is
assigned to its nearest unsubducted trench vertex $b$ of $p$, and the slab pull
per unit length applied at $b$ is
\begin{equation}
  f^{sp}_i = \frac{e_{sp}}{\Delta\ell_b}\sum_{v\to b} W_v\; n_i,
\end{equation}
where $\Delta\ell_b$ is the boundary length associated with $b$, $n_i$ the
outward normal of $p$ (perpendicular to the trench), and $e_{sp}$ the fraction of
slab weight transmitted to the surface plate, the remainder being supported by
viscous resistance of the mantle \cite{ConradLithgow2002}.

\paragraph{Interface friction.}
The plate interface extends downdip to depth $z_{\mathrm{int}}$. Its strength is
computed from the slab columns within map distance
$d_{\mathrm{int}} = z_{\mathrm{int}}/\tan\theta_d$ of the trench: their
area-weighted mean damage $\bar D$ and hydration $\bar w$. Integrating frictional
strength along the slanted interface, with overburden from the overriding
column,
\begin{equation}
  F_{\mathrm{int}} = \frac{\mu_f\,(1-a_b\bar D)\,\big(1-\lambda(\bar w)\big)}{\sin\theta_d}
  \int_0^{z_{\mathrm{int}}} P_s^{(q)}(\zeta)\dd\zeta.
\end{equation}
A wet, damaged interface is therefore weak, as megathrusts are observed to be.
Its horizontal projection acts on both plates, opposing their relative motion:
\begin{equation}
  f^{\mathrm{int}}_i = -F_{\mathrm{int}}\cos\theta_d\,\widehat{\Delta u}_i
  \ \text{on } p,
  \qquad
  +F_{\mathrm{int}}\cos\theta_d\,\widehat{\Delta u}_i
  \ \text{on } q.
\end{equation}

\paragraph{Trench suction.}
Flow around the sinking slab draws the overriding plate toward the trench. This
is represented by a traction on the trench of $q$, directed outward from $q$,
\begin{equation}
  f^{\mathrm{suc}}_i = \xi_s\, f^{sp}\, n^{(q)}_i,
\end{equation}
where $f^{sp}$ is the magnitude of the slab pull at the corresponding point of
$p$. It allows trench retreat and extension of the overriding plate.

\paragraph{Arc processes.}
Overriding columns of $q$ within map distance $d_{\mathrm{arc}}$ of the trench
receive hydration released from the slab and crust added by arc magmatism, both
in proportion to the convergence rate $v_c = \mac{v_n}$:
\begin{equation}
  S_w = k_w\,\bar w_{\mathrm{int}}\,\frac{v_c}{d_{\mathrm{arc}}},
  \qquad
  \frac{1}{A_v}\frac{\mathrm{D}V_c}{\mathrm{D}t}\bigg|_{\mathrm{arc}}
  = q_{\mathrm{arc}}\,\frac{v_c}{d_{\mathrm{arc}}},
\end{equation}
where $\bar w_{\mathrm{int}}$ is the mean hydration of slab columns at the
interface, $k_w$ a dimensionless hydration efficiency, and $q_{\mathrm{arc}}$
the thickness of crust produced per unit length of plate subducted. Arc crust
has reference density $\rho_{\mathrm{arc}}$ and is combined by the mixing rule.

\paragraph{Termination by collision.}
If the unsubducted trench columns of $p$ acquire buoyancy
$B \ge B_{\mathrm{crit}}$, for example because scraped-off crust accumulates
there or a continent or plateau arrives, subduction ceases along that stretch.
The slab columns assigned to it are deleted (slab break-off), and the plates are
welded as in Section~\ref{sec:contact}.

%----------------------------------------------------------------------
\subsection{Rifting and plate splitting}
\label{sec:rifting}
%----------------------------------------------------------------------

The mantle model supplies an upwelling rate $\Upsilon(\mathbf{x},t)$ at the base
of the plates. Upwelling in excess of a threshold $\Upsilon_c$ thins the plate
from below and, where the plate is in tension, damages it:
\begin{equation}
  S_L^{\Upsilon} = -\mac{\Upsilon - \Upsilon_c},
  \qquad
  S_D^{\Upsilon} = k_{\Upsilon}\,\frac{\mac{\Upsilon-\Upsilon_c}}{L}\,\mac{\chi}\,(1-D).
\end{equation}
Thinning lowers the strength of the column through $G$; damage lowers it
directly. Whether a rift develops is then decided by the column rheology and the
force balance, not prescribed.

A plate is split when a connected path of columns with $D \ge D_{\mathrm{split}}$
either joins two points on the plate boundary or closes on itself. Vertices
along the path are duplicated and the mesh is divided into two plates, whose
new common boundary is then classified as in
Section~\ref{sec:classification}. Splitting and welding are the only operations
that change the number of plates.

%----------------------------------------------------------------------
\subsection{Coverage}
%----------------------------------------------------------------------

The configuration must satisfy a coverage invariant: every point of the sphere
lies in exactly one unsubducted column, except within trenches, where subducted
columns of one plate lie beneath unsubducted columns of another. Gaps narrower
than $\Delta x_{\mathrm{acc}}$ and overlaps resisted by contact forces are
transient departures. The invariant is checked each step with the spatial index;
an unresolvable violation indicates a numerical error.

%======================================================================
\section{Force balance}
\label{sec:forcebalance}
%======================================================================

With lithostatic vertical stress and incompressibility, the depth-integrated
total horizontal stress of a column is
\begin{equation}
  \NN_{ij} = T_{ij} + \delta_{ij}T_{kk} - \delta_{ij}\Gamma.
\end{equation}
In the interior of each plate,
\begin{equation}
  \partial_j\NN_{ij} + t_i = 0,
  \label{eq:balance}
\end{equation}
where $t_i$ is the traction exerted on the plate base by the underlying mantle,
modelled as linear drag toward the mantle velocity $u^m_i$ supplied by the
mantle model:
\begin{equation}
  t_i = -C_b\left(u_i - u^m_i\right),
  \qquad
  C_b = \frac{\eta_a}{h_a}\left(\frac{L}{L_{\mathrm{ref}}}\right)^{m_b}.
  \label{eq:drag}
\end{equation}
Here $\eta_a$ is the asthenospheric viscosity, $h_a$ the thickness of the basal
shear zone, and the factor $(L/L_{\mathrm{ref}})^{m_b}$ represents the stronger
coupling of thick plates whose base reaches faster-moving mantle.

On the boundary of each plate,
\begin{equation}
  \NN_{ij}n_j = -\Gamma_{\mathrm{out}}\,n_i + f^{\partial}_i,
  \label{eq:bc}
\end{equation}
where $\Gamma_{\mathrm{out}}$ is the gravitational potential energy of the
material across the boundary (the neighbouring plate's boundary column at
contacts and trenches, $\Gamma_0$ at a gap) and $f^{\partial}_i$ is the sum of
the boundary forces that apply at that point: contact normal force and friction
(Section~\ref{sec:contact}), slab pull and interface friction on a subducting
plate, and interface friction and trench suction on an overriding plate
(Section~\ref{sec:convergent}). Because boundary forces depend on the velocities
of both plates, the force balances of all plates are solved together.

%======================================================================
\section{Evolution of the state}
\label{sec:kinematics}
%======================================================================

For unsubducted columns:
\begin{align}
  \frac{\mathrm{D}V_c}{\mathrm{D}t} &= \text{transfers at trenches and arc addition only},\\
  \frac{\mathrm{D}L}{\mathrm{D}t} &= -L\,\dot\varepsilon_{kk} + S_L + S_L^{\Upsilon},\\
  \frac{\mathrm{D}\rho_{0,c}}{\mathrm{D}t} &= \frac{\mathrm{D}\rho_{0,m}}{\mathrm{D}t} = 0
  \quad\text{(except through the mixing rule)},\\
  \frac{\mathrm{D}w}{\mathrm{D}t} &= S_w,
\end{align}
with $D$ and $T_{ij}$ governed by Eqs.~\eqref{eq:damage} and~\eqref{eq:evp}, and
$w$, $D$ clamped to $[0,1]$. Crustal thickness changes through $A_v$ as the mesh
deforms. The thermal source $S_L$ represents conductive cooling toward a
maximum thickness, with relaxation of mechanically over-thickened plate:
\begin{equation}
  S_L =
  \begin{cases}
    \dfrac{a_L^2\kappa}{2L}, & L < L_{\max},\\[2ex]
    -\dfrac{L - L_{\max}}{\tau_L}, & L \ge L_{\max},
  \end{cases}
  \qquad
  L_{\max} = L_{\max}^{o} + \left(L_{\max}^{cr}-L_{\max}^{o}\right)
  \operatorname{clamp}\!\left(\frac{\rho_{a,0}-\rho_{0,m}}{\rho_{a,0}-\rho_{\mathrm{dep}}},0,1\right),
\end{equation}
where $\kappa$ is thermal diffusivity, $a_L$ the coefficient of half-space
cooling ($L = a_L\sqrt{\kappa t}$ at age $t$), $L_{\max}^{o}$ and
$L_{\max}^{cr}$ the maximum thicknesses of oceanic and cratonic plates, and
$\rho_{\mathrm{dep}}$ the reference density of strongly depleted mantle. The
dependence of $L_{\max}$ on $\rho_{0,m}$ allows depleted mantle to sustain thick
cratonic roots. Subducted columns have frozen state.

%======================================================================
\section{Solution strategy}
\label{sec:solution}
%======================================================================

\paragraph{Rigid-plate limit.}
When all plates are much stronger than the forces acting on them, each plate
moves as a rigid rotation, $u_i = (\boldsymbol\omega_p\times\mathbf{r})_i$, with
angular velocity $\boldsymbol\omega_p$. Integrating Eq.~\eqref{eq:balance} over
the plate and using Eq.~\eqref{eq:bc} eliminates the unknown internal stress and
gives three torque-balance equations per plate,
\begin{equation}
  \oint_{\partial\Omega_p}\mathbf{r}\times\left(-\Gamma_{\mathrm{out}}\mathbf{n}
  + \mathbf{f}^{\partial}\right)\dd\ell
  + \int_{\Omega_p}\mathbf{r}\times\mathbf{t}\dd A = \mathbf{0},
  \label{eq:torque}
\end{equation}
coupled between plates through the boundary forces. This is the classical
plate-driving-force problem \cite{ForsythUyeda1975}, solved here with forces
computed from the column state. Friction makes the system mildly nonlinear; it
is solved by Picard or Newton iteration.

\paragraph{Deformable plates.}
The general problem solves Eqs.~\eqref{eq:evp}, \eqref{eq:balance}, and
\eqref{eq:bc} for the velocity field within all plates simultaneously. The
rigid-plate solution provides a starting guess and a check: plates whose
interior stress resultants remain below yield should deform negligibly.

\paragraph{Time step.}
Each step proceeds in the following order:
\begin{enumerate}
  \item Reconstruct profiles and column properties from the state
        (Sections~\ref{sec:thermal}--\ref{sec:integrated}).
  \item Compute contact geometry and boundary forces
        (Section~\ref{sec:boundaries}).
  \item Solve for velocities (rigid or deformable).
  \item Advect vertices; move slab columns with their trenches.
  \item Apply boundary events: create plate in gaps, flag newly subducted
        columns and scrape their crust, delete slab columns beyond
        $d_{\max}$, initiate subduction, weld, and split.
  \item Update $D$, $L$, $w$, $V_c$, and $T_{ij}$.
  \item Split long edges and remesh where triangles have degraded, remapping
        extensive quantities ($V_c$, and area-weighted intensive ones)
        conservatively.
  \item Check coverage.
\end{enumerate}

%======================================================================
\section{Parameters}
%======================================================================

Status \textsf{T} denotes a typical literature value; \textsf{P} a placeholder
to be calibrated; \textsf{N} a numerical parameter whose value should not affect
results when chosen sensibly. Calibration targets include mean elevations of
ridge crests, old seafloor, and continents; seismogenic depths; effective
elastic thicknesses; present-day plate speeds and their correlation with
attached slab length; and the requirement that undamaged lithosphere resist the
available driving forces (of order \SI{e12}{} to \SI{e13}{\newton\per\metre})
while damaged lithosphere does not.

\begin{longtable}{@{}lp{7.0cm}rlc@{}}
\caption{Global parameters.}\label{tab:params}\\
\toprule
Symbol & Meaning & Nominal & Units & Status\\
\midrule
\endfirsthead
\toprule
Symbol & Meaning & Nominal & Units & Status\\
\midrule
\endhead
\bottomrule
\endfoot
\multicolumn{5}{@{}l}{\emph{Thermal}}\\
$T_s$ & surface temperature & 0 & \si{\celsius} & T\\
$T_L$ & temperature defining the plate base & 1300 & \si{\celsius} & T\\
$T_p$ & mantle potential temperature & 1350 & \si{\celsius} & T\\
$\gamma_{\mathrm{ad}}$ & adiabatic gradient & 0.4 & \si{\kelvin\per\kilo\metre} & T\\
$\kappa$ & thermal diffusivity & \num{1e-6} & \si{\square\metre\per\second} & T\\
$a_L$ & half-space plate-thickness coefficient & 2.32 & -- & T\\
$L_{\max}^{o}$ & maximum oceanic plate thickness & 110 & \si{\kilo\metre} & T\\
$L_{\max}^{cr}$ & maximum cratonic plate thickness & 230 & \si{\kilo\metre} & T\\
$\tau_L$ & relaxation time of over-thickened plate & 50 & Myr & P\\
\midrule
\multicolumn{5}{@{}l}{\emph{Density and isostasy}}\\
$g$ & gravitational acceleration & 9.81 & \si{\metre\per\square\second} & T\\
$T_0$ & reference temperature for reference densities & 0 & \si{\celsius} & --\\
$\alpha$ & thermal expansion coefficient & \num{3.2e-5} & \si{\per\kelvin} & T\\
$\rho_w$ & seawater density & 1030 & \si{\kilo\gram\per\cubic\metre} & T\\
$\rho_{a,0}$ & reference density of asthenosphere & 3380 & \si{\kilo\gram\per\cubic\metre} & T\\
$\rho_{\mathrm{dep}}$ & reference density of strongly depleted mantle & 3310 & \si{\kilo\gram\per\cubic\metre} & T\\
$z_c$ & compensation depth & 300 & \si{\kilo\metre} & --\\
$e_{\mathrm{ref}}$ & elevation of the reference column & $-2.5$ & \si{\kilo\metre} & T\\
\midrule
\multicolumn{5}{@{}l}{\emph{Friction and pore pressure}}\\
$\mu_f$ & friction coefficient & 0.6 & -- & T\\
$\lambda_0$ & pore-pressure ratio, dry (hydrostatic) & 0.37 & -- & T\\
$\lambda_1$ & pore-pressure ratio, fully hydrated & 0.9 & -- & P\\
\midrule
\multicolumn{5}{@{}l}{\emph{Creep and plasticity}}\\
$n$ & creep stress exponent & 3.5 & -- & T\\
$R_g$ & gas constant & 8.314 & \si{\joule\per\mole\per\kelvin} & --\\
$\dot\varepsilon_{\mathrm{ref}}$ & reference strain rate & \num{1e-15} & \si{\per\second} & T\\
$\rho^{(a)}, T_{bd}^{(a)}, Q^{(a)}$ & closure anchors & Table~\ref{tab:anchors} & & T/P\\
$\Delta T_w$ & reduction of $T_{bd}$ at full hydration & 150 & \si{\kelvin} & P\\
$\sigma_{\max}$ & low-temperature plasticity cap & 800 & \si{\mega\pascal} & P\\
$\kappa_\sigma$ & differential-stress-to-invariant factor & 0.5 & -- & P\\
$\sigma_{\mathrm{thr}}$ & strength threshold for elastic layers & 10 & \si{\mega\pascal} & P\\
\midrule
\multicolumn{5}{@{}l}{\emph{Elasticity}}\\
$G_s$ & shear modulus & 30 & \si{\giga\pascal} & T\\
$E$ & Young's modulus (flexure) & 70 & \si{\giga\pascal} & T\\
$\nu$ & Poisson's ratio (flexure) & 0.25 & -- & T\\
\midrule
\multicolumn{5}{@{}l}{\emph{Damage}}\\
$k_D$ & damage per unit viscoplastic strain & 10 & -- & P\\
$a_b$ & frictional weakening at full damage & 0.7 & -- & P\\
$a_d$ & creep weakening at full damage & 0.9 & -- & P\\
$\tau_{h,0}$ & healing time at $T_{h,0}$ & 100 & Myr & P\\
$T_{h,0}$ & reference healing temperature & 400 & \si{\celsius} & P\\
$Q_h$ & healing activation energy & 300 & \si{\kilo\joule\per\mole} & P\\
$\zeta_h$ & depth at which healing temperature is evaluated & 15 & \si{\kilo\metre} & P\\
$\ell_D$ & shear-zone width scale & 15 & \si{\kilo\metre} & P\\
\midrule
\multicolumn{5}{@{}l}{\emph{Basal coupling}}\\
$\eta_a$ & asthenospheric viscosity & \num{1e20} & \si{\pascal\second} & T\\
$h_a$ & basal shear-zone thickness & 100 & \si{\kilo\metre} & P\\
$L_{\mathrm{ref}}$ & reference plate thickness for coupling & 100 & \si{\kilo\metre} & P\\
$m_b$ & exponent of thickness dependence of coupling & 1 & -- & P\\
\midrule
\multicolumn{5}{@{}l}{\emph{New plate}}\\
$L_0$ & initial plate thickness & 5 & \si{\kilo\metre} & P\\
$h_c^{o}$ & oceanic crustal thickness & 7 & \si{\kilo\metre} & T\\
$\rho_{0,c}^{o}$ & oceanic crust reference density & 2950 & \si{\kilo\gram\per\cubic\metre} & T\\
$\rho_{0,m}^{o}$ & oceanic mantle reference density & 3360 & \si{\kilo\gram\per\cubic\metre} & T\\
$w^{o}$ & hydration of new plate & 0 & -- & P\\
\midrule
\multicolumn{5}{@{}l}{\emph{Subduction}}\\
$B_{\mathrm{crit}}$ & buoyancy above which a column cannot subduct & \num{5e6} & \si{\kilo\gram\per\square\metre} & P\\
$\rho_{\mathrm{scr}}$ & crustal density below which crust is scraped off & 2850 & \si{\kilo\gram\per\cubic\metre} & P\\
$f_{\mathrm{acc}}$ & fraction of scraped crust accreted to the overriding plate & 0.5 & -- & P\\
$\theta_d$ & representative slab dip & 45 & \si{\degree} & T\\
$z_{\mathrm{slab}}$ & depth limit of slab pull & 660 & \si{\kilo\metre} & T\\
$\rho_{\mathrm{ecl}}$ & reference density of eclogitized slab crust & 3500 & \si{\kilo\gram\per\cubic\metre} & T\\
$e_{sp}$ & fraction of slab weight transmitted to the plate & 0.2 & -- & P\\
$z_{\mathrm{int}}$ & depth extent of the plate interface & 40 & \si{\kilo\metre} & T\\
$\xi_s$ & trench suction relative to slab pull & 0.1 & -- & P\\
$d_{\mathrm{arc}}$ & width of the arc region & 150 & \si{\kilo\metre} & T\\
$k_w$ & hydration efficiency of slab fluids & 1 & -- & P\\
$q_{\mathrm{arc}}$ & arc crust produced per unit length subducted & 1 & \si{\kilo\metre} & P\\
$\rho_{\mathrm{arc}}$ & reference density of arc crust & 2800 & \si{\kilo\gram\per\cubic\metre} & T\\
\midrule
\multicolumn{5}{@{}l}{\emph{Rifting and splitting}}\\
$\Upsilon_c$ & upwelling threshold for basal thinning & -- & \si{\metre\per\second} & P\\
$k_\Upsilon$ & damage efficiency of upwelling & 1 & -- & P\\
$D_{\mathrm{split}}$ & damage threshold for plate splitting & 0.9 & -- & P\\
\midrule
\multicolumn{5}{@{}l}{\emph{Numerical}}\\
$g_{\mathrm{tol}}$ & contact tolerance & \textasciitilde mesh/10 & \si{\metre} & N\\
$\Delta x_{\mathrm{acc}}$ & gap width that triggers plate creation & \textasciitilde mesh/2 & \si{\metre} & N\\
$\ell_{\max}$ & maximum boundary edge length & \textasciitilde mesh & \si{\metre} & N\\
$k_c$ & contact penalty stiffness & large & \si{\newton\per\square\metre} & N\\
$v_{\mathrm{reg}}$ & friction regularization speed & 0.1 & mm/yr & N\\
\end{longtable}

%======================================================================
\section{Assumptions and limitations}
%======================================================================

\begin{itemize}
  \item \textbf{Thin-sheet kinematics.} Horizontal velocity is uniform with depth
        in each column; decoupling of crust from mantle is represented only
        through the effective elastic thickness.
  \item \textbf{Linear geotherm.} Crustal radiogenic heating and the
        error-function profile of young oceanic plate are neglected; resulting
        errors are partly absorbed by calibration.
  \item \textbf{Phenomenological closures.} The density-based choice of creep
        properties, the hydration and damage weakening laws, and the rules for
        scraping, arc production, and slab pull transmission reproduce
        qualitative behaviour and ordering, not laboratory or field
        quantities.
  \item \textbf{Absolute densities.} An error of
        \SI{10}{\kilo\gram\per\cubic\metre} in a reference density over
        \SI{200}{\kilo\metre} of column changes isostatic elevation by about
        \SI{0.6}{\kilo\metre}. Reference densities should be calibrated against
        elevations.
  \item \textbf{Slab representation.} The slab has a single representative dip
        and is passive: it transmits a fixed fraction of its weight and has no
        internal dynamics, thermal assimilation, or interaction with the
        mantle flow beyond its removal at $d_{\max}$.
  \item \textbf{Mantle coupling is one-way.} The mantle model supplies $u^m_i$
        and $\Upsilon$; plates and slabs do not feed back on mantle flow.
  \item \textbf{Fixed brittle--ductile partition} at the reference strain rate,
        and incompressible in-plane elasticity alongside a compressible
        Poisson's ratio for flexure.
  \item \textbf{Topology changes} (creation, subduction, welding, splitting)
        are discrete events applied between force-balance solutions; their
        rules involve thresholds whose influence on results should be tested.
\end{itemize}

%======================================================================
\begin{thebibliography}{99}
%======================================================================

\bibitem{BercoviciRicard2012}
D.~Bercovici and Y.~Ricard,
Mechanisms for the generation of plate tectonics by two-phase grain-damage and
pinning,
\emph{Physics of the Earth and Planetary Interiors} \textbf{202--203}, 27--55 (2012).

\bibitem{BurovDiament1995}
E.~B. Burov and M.~Diament,
The effective elastic thickness ($T_e$) of continental lithosphere: What does it
really mean?
\emph{Journal of Geophysical Research} \textbf{100}(B3), 3905--3927 (1995).

\bibitem{Byerlee1978}
J.~Byerlee,
Friction of rocks,
\emph{Pure and Applied Geophysics} \textbf{116}, 615--626 (1978).

\bibitem{ConradLithgow2002}
C.~P. Conrad and C.~Lithgow-Bertelloni,
How mantle slabs drive plate tectonics,
\emph{Science} \textbf{298}, 207--209 (2002).

\bibitem{EnglandMcKenzie1982}
P.~England and D.~McKenzie,
A thin viscous sheet model for continental deformation,
\emph{Geophysical Journal of the Royal Astronomical Society} \textbf{70},
295--321 (1982).

\bibitem{ForsythUyeda1975}
D.~Forsyth and S.~Uyeda,
On the relative importance of the driving forces of plate motion,
\emph{Geophysical Journal of the Royal Astronomical Society} \textbf{43},
163--200 (1975).

\bibitem{GleasonTullis1995}
G.~C. Gleason and J.~Tullis,
A flow law for dislocation creep of quartz aggregates determined with the molten
salt cell,
\emph{Tectonophysics} \textbf{247}, 1--23 (1995).

\bibitem{HirthKohlstedt2003}
G.~Hirth and D.~Kohlstedt,
Rheology of the upper mantle and the mantle wedge: A view from the
experimentalists,
in \emph{Inside the Subduction Factory}, Geophysical Monograph 138, American
Geophysical Union, 83--105 (2003).

\bibitem{Mackwell1998}
S.~J. Mackwell, M.~E. Zimmerman, and D.~L. Kohlstedt,
High-temperature deformation of dry diabase with application to tectonics on
Venus,
\emph{Journal of Geophysical Research} \textbf{103}(B1), 975--984 (1998).

\bibitem{Saramito2009}
P.~Saramito,
A new elastoviscoplastic model based on the Herschel--Bulkley viscoplastic
model,
\emph{Journal of Non-Newtonian Fluid Mechanics} \textbf{158}, 154--161 (2009).

\end{thebibliography}

\end{document}
